% Секция 1.2 — Современные представления о механизмах нейронного кодирования информации в мозге
% Драфт для ревью. Не включать напрямую — после одобрения перенести в part1.tex.

\section{Современные представления о механизмах нейронного кодирования информации в мозге}
\label{sec:coding}

Мозг непрерывно решает задачу извлечения релевантной информации из сложной и многомерной среды. Одним из центральных вопросов нейронауки является вопрос о \textit{нейронном коде}~--- способе, которым нервные клетки представляют информацию о внешнем мире и внутренних состояниях организма~\cite{DayanAbbott2001}. Исторически в изучении нейронного кода сложились два основных подхода: анализ активности отдельных клеток (индивидуальный код) и исследование коллективных паттернов популяции нейронов (популяционный код). Как будет показано ниже, ни один из этих подходов в отдельности не даёт исчерпывающего описания механизмов кодирования, что привело к развитию представлений о смешанной селективности нейронов.

\subsection{Индивидуальный код}
\label{sec:local_coding}

Наиболее наглядным примером индивидуального нейронного кода служит \textit{локальное кодирование}, при котором каждое уникальное событие или контекст ассоциируется с активностью одной клетки. Классической демонстрацией этого принципа стали эксперименты Хьюбела и Визеля~\cite{HubelWiesel1962}, показавшие, что нейроны первичной зрительной коры избирательно реагируют на ориентацию краёв в зрительном поле. Скорость генерации потенциалов действия такого нейрона может быть описана как функция одной переменной~--- ориентации стимула, а не интенсивности всех пикселей изображения.

Ярким примером локального кодирования на высшем когнитивном уровне являются так называемые \textit{клетки концепта} (concept cells). Было показано, что нейроны медиальной височной доли человека избирательно реагируют на конкретных людей или объекты~--- таких как Дженнифер Энистон, Люк Скайуокер или Пизанская башня~--- независимо от модальности предъявления стимула (фотография, текст, звук)~\cite{concepts}. Активация такого нейрона отражает распознавание абстрактного концепта соответствующего объекта, а не его конкретные сенсорные признаки.

Важнейшим классом нейронов с выраженной индивидуальной селективностью являются \textit{клетки места} гиппокампа. О'Киф и Достровский обнаружили, что отдельные нейроны зоны CA1 гиппокампа крысы избирательно активируются, когда животное находится в определённой области пространства, образуя так называемые \textit{поля места}~\cite{o1971hippocampus}. Комбинация полей различных клеток места формирует внутреннюю пространственную карту среды. При этом местоположение нейронов в мозге не отражает топографию их полей: расположение нейронов нетопографично, хотя размер полей увеличивается от дорсального к вентральному гиппокампу.

Помимо клеток места, к системе пространственного кодирования относятся \textit{координатные клетки} (grid cells) медиальной энторинальной коры, образующие регулярную гексагональную решётку полей, покрывающую всё доступное животному пространство~\cite{Hafting2005}, а также \textit{клетки направления головы}~\cite{Taube1990}. Совместно с клетками скорости~\cite{kropff2015speed} и клетками границ эти типы нейронов формируют многокомпонентную систему обработки пространственной информации в мозге.

Наряду с пространственными переменными, в гиппокампе были обнаружены нейроны, кодирующие время. \textit{Клетки времени} (time cells) последовательно активируются в определённые моменты временного интервала, формируя распределённое представление времени, что было показано как у грызунов, так и у человека~\cite{Umbach2020, Cao2022}.

Несмотря на впечатляющие примеры узкоспециализированных нейронов, чисто локальное кодирование имеет существенные ограничения. Локальные коды страдают от низкой репрезентативной способности: совокупность из $N$ нейронов способна кодировать не более $N$ различных контекстов. Кроме того, локальные коды требуют большого количества нейронов для охвата всего пространства возможных стимулов. Эти теоретические ограничения, а также экспериментальные свидетельства о более сложных паттернах нейронной активности привели к развитию популяционного подхода.

\subsection{Популяционный код}
\label{sec:population_coding}

Противоположностью локального является \textit{плотное кодирование}, при котором каждый контекст представлен объединённой активностью всех нейронов популяции. Теоретически плотное кодирование обеспечивает высокую репрезентативную способность: при дискретизации активности на $M$ уровней $N$ нейронов позволяют кодировать порядка $M^N$ контекстов. Однако плотное кодирование неизбежно страдает от перекрёстного взаимодействия, поскольку каждая клетка вовлечена в представление каждого контекста. Компромиссным решением служит \textit{разреженное кодирование}, при котором каждый контекст кодируется различным поднабором нейронов, что снижает общую нейронную активность, необходимую для представления информации.

Большинство сенсорных, когнитивных и моторных функций зависят от взаимодействия многих нейронов~\cite{Averbeck2006}. Хотя спайки отдельных клеток стохастичны и нелинейно зависят от входных сигналов, на уровне популяции формируются устойчивые паттерны совместной активности. Это наблюдение послужило основой для перехода от изучения отдельных нейронов к анализу популяционной активности~\cite{Pouget2000}.

Активность популяции из $N$ одновременно записанных нейронов может быть описана в терминах $N$-мерного пространства состояний, где каждая координата соответствует активности отдельной клетки. В экспериментах было показано, что реальная активность занимает лишь малую часть этого пространства и сосредоточена вблизи низкоразмерного подпространства~--- \textit{нейронного многообразия} (neural manifold)~\cite{Gallego2017}. Это многообразие определяется набором базисных векторов, которые получили название \textit{нейронных мод}. Активность каждого нейрона представляет собой взвешенную комбинацию нейронных мод, веса которой меняются во времени, задавая \textit{латентную динамику} популяции~\cite{Cunningham2014}.

Форма и структура нейронного многообразия во многом определяются физическими связями между нейронами сети. Было показано, что нейронное подпространство слабо подвержено изменениям при обучении на масштабах порядка часов~\cite{Gallego2018}. При этом многообразия нейронной активности по своей природе нелинейны~\cite{De2023, Fortunato2023}, и их низкая размерность часто скрывается методами анализа, предполагающими линейность данных.

Таким образом, популяционный подход позволяет описывать коллективную динамику нейронов с помощью небольшого числа латентных переменных. Это существенно расширяет возможности анализа по сравнению с индивидуальным описанием, однако ставит новые вопросы: каков вклад отдельного нейрона в коллективную активность и какие именно переменные внешнего мира кодирует каждая клетка?

\subsection{Смешанная селективность}
\label{sec:mixed_selectivity}

Исследования последних лет показали, что разделение на индивидуальный и популяционный подходы существенно упрощает реальную картину нейронного кодирования. Явление \textit{смешанной селективности} (mixed selectivity), при котором отдельный нейрон отражает комбинации нескольких стимулов, поведенческих переменных или когнитивных состояний, а не единственный признак, оказалось широко распространённым в мозге~\cite{TyeMiller2024}. Нейроны префронтальной коры, например, приобретают смешанную селективность к релевантным переменным задачи в процессе обучения~\cite{mix3, Mante2013}. Более того, даже канонические примеры функционально специализированных нейронов кодируют множество переменных: клетки места отражают не только положение животного, но и время, награду и траекторию движения~\cite{Zheng2021, Lian2022}. Аналогичный принцип проявляется и в сенсорной обработке: нейроны первичной зрительной коры и ретроспленальной коры способны кодировать как текущие, так и запомненные зрительные образы~\cite{Makino2015}.

Смешанная селективность создаёт фундаментальную \textit{проблему атрибуции}: кажущаяся селективность к одной переменной может объясняться селективностью к коррелированному ковариату. Эта проблема особенно остро стоит в натуралистических экспериментальных условиях, где поведенческие переменные внутренне коррелированы. Например, активность нейрона может выглядеть как кодирование локомоции, тогда как на самом деле она определяется скоростью; а клетка места может быть ошибочно интерпретирована как реагирующая на социальное взаимодействие лишь потому, что контакты между животными приурочены к определённым участкам арены.

Эта проблема проявляется в нейронауке повсеместно. Мусалл и коллеги~\cite{Musall2019} показали, что неинструктированные движения (uninstructed movements) объясняют значительную часть вариации нейронной активности по всей коре головного мозга, а значит, кажущаяся когнитивная селективность корковых нейронов может быть артефактом неконтролируемых моторных переменных. В гиппокампе тета-модуляция активности осложняет разграничение кодирования скорости и фазы~\cite{Vanderwolf1969, McNaughton1983}. Эти примеры указывают на общее требование к методам анализа нейронной селективности: они должны явным образом учитывать коррелированные ковариаты и временную структуру нейронных и поведенческих сигналов.

Развитие технологий автоматического анализа поведения, включая методы неконтролируемого~\cite{Weinreb2024}, контролируемого~\cite{Harris2023deepaction, Segalin2021} и основанного на правилах~\cite{Gabriel2022} распознавания, позволяет извлекать из видеозаписей десятки и сотни поведенческих признаков. Это создаёт богатые, но высококоррелированные наборы данных, усугубляя проблему атрибуции: при наличии множества взаимозависимых переменных установить, какая именно из них определяет активность конкретного нейрона, становится особенно сложно.

Таким образом, смешанная селективность является не исключением, а правилом нейронного кодирования. Это обстоятельство требует разработки аналитических подходов, которые способны не только обнаруживать связь между активностью нейрона и внешними переменными, но и разделять вклады коррелированных факторов, отличая истинное кодирование от ложных ассоциаций, обусловленных ковариацией поведенческих переменных.
